\section{Related work}
\label{sec:related}

\paragraph{The 2009 attempt.} Bailey et al.~\cite{bailey2009} fixed
the iteration function this campaign walks, $R\mapsto\sigma^j(R)+R$
with $j$ from the weight of $x$, the weight-34 distinguished-point
rule, and the $2^{60.9}$ budget, and reported per-platform rates on
Core~2 (533 cycles per iteration with 128-bit vectors), Cell, GPUs,
FPGAs~\cite{fan2010fpga} and ASIC estimates; their SHARCS
paper~\cite{bailey2009sharcs} priced the whole Certicom ECC2 list.
Bos, Kleinjung, Niederhagen and Schwabe's Cell
implementation~\cite{bos2010cell} is the bitsliced lineage this
client's CPU and first GPU backends belong to. The type-II
optimal-polynomial-basis multiplication of Bernstein and
Lange~\cite{bernsteinlange2010onb} is used verbatim; the conversion
networks generated here match their hand-designed chains to within a
few instructions (Section~\ref{sec:arith-bitsliced}). Bernstein's bit
operation counts for binary multiplication~\cite{bernstein2009batch}
are the reference the three-input-logic finding is measured against.
Bernstein et al.'s later FPGA designs~\cite{bernstein2016fpga}, which
set the current size record for a completed binary-curve computation
by solving a $2^{117.35}$-order instance on sect113r2 and were
extrapolated to a 127-bit field, are the calibration for the pre-board
estimate of Section~\ref{sec:fpga}, which the measurement then met.
Harley's group solved ECC2K-95 in 1998 and ECC2K-108 in April 2000,
the latter with about $1{,}300$ volunteers' machines over four
months~\cite{harley1998,harley2000}; those two are the only Koblitz
challenges completed, and their walks are the ancestors of the one in
Section~\ref{sec:walk-sigma}.

\paragraph{Rho on automorphism classes.} The $\sqrt{2m}$ saving from
Frobenius and negation on Koblitz curves is due to Wiener and
Zuccherato~\cite{wienerzuccherato1998} and Gallant, Lambert and
Vanstone~\cite{glv2000}, generalised by Duursma, Gaudry and
Morain~\cite{duursma1999}. Teske's $r$-adding walks~\cite{teske2001}
and the Bernstein--Lange--Schwabe treatment of fruitless cycles under
the negation map~\cite{bls2011} are the background to
Section~\ref{sec:walk-table}; the cycle classes measured there are the
$\tau$-relation analogue, and the measured walk constants of
Table~\ref{tab:walkconst} are, to our knowledge, the first reported
for this branch distribution on this family. The distinguished-point
method is van Oorschot and Wiener's~\cite{vanoorschot1999}. The 112-bit
prime-field record of Bos et al.~\cite{bos2012ps3} is the nearest
completed computation in scale.

\paragraph{Index calculus on binary curves.} Semaev's summation
polynomials~\cite{semaev2004}, Gaudry's and Diem's index
calculus~\cite{gaudry2009,diem2011}, the Faug\`ere--Perret--Petit--Renault
and Petit--Quisquater analyses of Weil-descent
systems~\cite{faugere2012,petitquisquater2012} and Galbraith and
Gebregiyorgis's characteristic-two
algorithms~\cite{galbraithgebregiyorgis2014} are the methods
Section~\ref{sec:alternatives} prices on this curve. The structural
closure at $n=131$ (no Frobenius-stable subspace of intermediate
dimension) is elementary once the normal basis is in view, and the
measured ladder confirms that the surviving weight-bounded bases do
not reach the budget.

\paragraph{Measurement discipline.} The companion
paper~\cite{ecdlppaper} states the boundary-unit-table-class rule and
applies it across rho, baby-step giant-step, kangaroo and prime-field
index calculus; this paper is its ECC2K-130 chapter at full length.
